\setcounter{section}{5}


\inputaufgabe
{}
{





Es sei $X$ ein
\definitionsverweis {topologischer Raum}{}{}
und ${ \mathcal  F   }$ eine
\definitionsverweis {Prägarbe}{}{}
auf $X$. Zeige, dass durch die Zuordnung

\mathdisp {U \mapsto \prod_{P \in U} { \mathcal  F   }_P} {  }

\zusatzklammer {die Produktmenge aus allen
\definitionsverweis {Halmen}{}{}
zu $U$} {} {}
mit den natürlichen Projektionen eine Prägarbe gegeben ist, und dass es einen natürlichen
\definitionsverweis {Prägarbenhomomorphismus}{}{}
von ${ \mathcal  F   }$  in diese Prägarbe gibt.





}
{} {}

\inputaufgabe
{}
{





Es sei ${ \mathcal  F   }$ eine
\definitionsverweis {Prägarbe}{}{}
auf einem
\definitionsverweis {topologischen Raum}{}{}
$X$ und $\widetilde{  { \mathcal  F   }  }$ die
\definitionsverweis {Vergarbung}{}{}
zu ${ \mathcal  F   }$. Zeige, dass es zu jedem
\definitionsverweis {Prägarben-Morphismus}{}{}
\maabbdisp {\psi} { { \mathcal  F   }  } { { \mathcal  G   }
} {}
in eine
\definitionsverweis {Garbe}{}{}
${ \mathcal  G   }$ eine eindeutige Faktorisierung
\maabbdisp {\tilde{\psi}} { \widetilde{  { \mathcal  F   }  } } { { \mathcal  G   }
} {}
gibt.





}
{} {}

Die Vergarbung einer konstanten Prägarbe nennt man \stichwort {lokal konstante Garbe} {} und manchmal auch einfach \stichwort {konstante Garbe} {.}

\inputaufgabe
{}
{





Es sei ${ \mathcal  F   }$ eine
\definitionsverweis {konstante Prägarbe}{}{}
auf einem
\definitionsverweis {topologischen Raum}{}{}
$X$ zur Menge $M$. Zeige, dass der
\definitionsverweis {Halm}{}{}
der
\definitionsverweis {Vergarbung}{}{}
von ${ \mathcal  F   }$ in jedem Punkt

\mavergleichskette
{\vergleichskette
{P
}
{ \in }{ X
}
{  }{
}
{    }{
}
{   }{
}
}
{}{}{}
gleich $M$ ist.





}
{} {}

\inputaufgabe
{}
{





Es sei $G$ eine
\definitionsverweis {diskrete}{}{}
\definitionsverweis {topologische Gruppe}{}{}
und $X$ ein
\definitionsverweis {topologischer Raum}{}{.}
Es sei ${ \mathcal  G   }$ die
\definitionsverweis {konstante Prägarbe}{}{}
auf $X$ zu $G$. Zeige, dass die
\definitionsverweis {Vergarbung}{}{}
von ${ \mathcal  G   }$ gleich
\mathl{C^0(-,G)}{} ist.





}
{} {}

\inputaufgabegibtloesung
{}
{





Es sei $X$ ein
\definitionsverweis {topologischer Raum}{}{,}
${ \mathcal  G   }$ eine
\definitionsverweis {Garbe}{}{}
auf $X$ und

\mavergleichskette
{\vergleichskette
{ { \mathcal  F   }
}
{ \subseteq }{ { \mathcal  G   }
}
{  }{
}
{    }{
}
{   }{
}
}
{}{}{}
eine
\definitionsverweis {Untergarbe}{}{.}
Es sei

\mavergleichskette
{\vergleichskette
{t
}
{ \in }{ \Gamma   \left(   X ,  { \mathcal  G   }   \right)
}
{  }{
}
{    }{
}
{   }{
}
}
{}{}{}
mit

\mavergleichskette
{\vergleichskette
{t_P
}
{ \in }{ { \mathcal  F   }_P
}
{  }{
}
{    }{
}
{   }{
}
}
{}{}{}
für alle

\mavergleichskette
{\vergleichskette
{P
}
{ \in }{ X
}
{  }{
}
{    }{
}
{   }{
}
}
{}{}{.}
Zeige

\mavergleichskette
{\vergleichskette
{t
}
{ \in }{ \Gamma   \left(   X ,  { \mathcal  F   }   \right)
}
{  }{
}
{    }{
}
{   }{
}
}
{}{}{.}





}
{} {}

\inputaufgabe
{}
{





Es sei $X$ ein
\definitionsverweis {topologischer Raum}{}{}
und es sei
\maabb {\varphi} { { \mathcal  F   } } { { \mathcal  G   }
} {}
ein
\definitionsverweis {Homomorphismus}{}{}
von
\definitionsverweis {Garben}{}{}
von
\definitionsverweis {kommutativen Gruppen}{}{.}
Zeige, dass durch

\mavergleichskette
{\vergleichskette
{ \left(  \operatorname{kern} \varphi  \right)  (U)
}
{ \defeq }{ \operatorname{kern} \varphi_U
}
{  }{
}
{    }{
}
{   }{
}
}
{}{}{}
eine Garbe von Gruppen auf $X$ definiert ist.





}
{} {}

\inputaufgabe
{}
{





Es sei $X$ ein
\definitionsverweis {topologischer Raum}{}{}
und es sei
\maabb {\varphi} { { \mathcal  F   } } { { \mathcal  G   }
} {}
ein
\definitionsverweis {Homomorphismus}{}{}
von
\definitionsverweis {Garben}{}{}
von
\definitionsverweis {kommutativen Gruppen}{}{.}
Zeige, dass $\varphi$ genau dann
\definitionsverweis {injektiv}{}{}
ist, wenn
\mathl{\operatorname{kern} \varphi}{} die Nullgarbe ist.





}
{} {}

\inputaufgabe
{}
{





Es sei $X$ ein
\definitionsverweis {topologischer Raum}{}{}
und es sei
\maabb {\varphi} { { \mathcal  F   } } { { \mathcal  G   }
} {}
ein
\definitionsverweis {Homomorphismus}{}{}
von
\definitionsverweis {Garben}{}{}
von
\definitionsverweis {kommutativen Gruppen}{}{.}
Zeige, dass $\varphi$ genau dann
\definitionsverweis {surjektiv}{}{}
ist, wenn

\mavergleichskette
{\vergleichskette
{ \operatorname{bild} \varphi
}
{ = }{ { \mathcal  G   }
}
{  }{
}
{    }{
}
{   }{
}
}
{}{}{}
ist.





}
{} {}

\inputaufgabe
{}
{





Es sei $X$ ein
\definitionsverweis {topologischer Raum}{}{}
und es sei
\maabb {\varphi} { { \mathcal  F   } } { { \mathcal  G   }
} {}
ein
\definitionsverweis {Homomorphismus}{}{}
von
\definitionsverweis {Garben}{}{}
von
\definitionsverweis {kommutativen Gruppen}{}{.}
Zeige, dass für jeden Punkt

\mavergleichskette
{\vergleichskette
{P
}
{ \in }{ X
}
{  }{
}
{    }{
}
{   }{
}
}
{}{}{}
die Beziehung

\mavergleichskette
{\vergleichskette
{ \left(  \operatorname{bild} \varphi  \right)_P
}
{ = }{ \operatorname{bild}  \left(  \varphi_P   \right)
}
{  }{
}
{    }{
}
{   }{
}
}
{}{}{}
ist.





}
{} {}

\inputaufgabe
{}
{





Es sei eine
\definitionsverweis {Garbe}{}{}
${ \mathcal  G   }$ von
\definitionsverweis {kommutativen Gruppen}{}{}
und eine
\definitionsverweis {Untergarbe}{}{}
von Gruppen

\mavergleichskette
{\vergleichskette
{ { \mathcal  F   }
}
{ \subseteq }{ { \mathcal  G   }
}
{  }{
}
{    }{
}
{   }{
}
}
{}{}{}
gegeben. Zeige, dass es einen kanonischen
\definitionsverweis {surjektiven}{}{}
\definitionsverweis {Garbenhomomorphismus von kommutativen Gruppen}{}{}
\maabbdisp {} { { \mathcal  G   } } { { \mathcal  G   } / { \mathcal  F   }
} {}
gibt.





}
{} {}

\inputaufgabe
{}
{





Es sei eine
\definitionsverweis {Garbe}{}{}
${ \mathcal  G   }$ von
\definitionsverweis {kommutativen Gruppen}{}{}
und eine
\definitionsverweis {Untergarbe}{}{}
von Gruppen

\mavergleichskette
{\vergleichskette
{ { \mathcal  F   }
}
{ \subseteq }{ { \mathcal  G   }
}
{  }{
}
{    }{
}
{   }{
}
}
{}{}{}
gegeben, und es sei
\mathl{{ \mathcal  G   } / { \mathcal  F   }}{} die
\definitionsverweis {Quotientengarbe}{}{.}
Zeige

\mavergleichskettedisp
{\vergleichskette
{ \left(  { \mathcal  G   } / { \mathcal  F   }   \right)_P
}
{ =} { { \mathcal  G   }_P / { \mathcal  F   }_P
}
{ } {
}
{ } {
}
{ } {
}
}
{}{}{}
für jeden Punkt

\mavergleichskette
{\vergleichskette
{P
}
{ \in }{ X
}
{  }{
}
{    }{
}
{   }{
}
}
{}{}{.}





}
{} {}
